% ---------------------------------------------------------------------
%  Pakete
% ---------------------------------------------------------------------
\usepackage[T1]{fontenc}
\usepackage[utf8]{inputenc}
\usepackage[ngerman]{babel}
\IfFileExists{FiraSans.sty}{%
  \usepackage[sfdefault,lining]{FiraSans}
  \usepackage[scale=0.85]{FiraMono}}{\usepackage{lmodern}}
\usepackage{amsmath,amssymb}
\usepackage[locale=DE]{siunitx}
\usepackage{braket}
\usepackage{booktabs}
\usepackage{listings}
\usepackage[most]{tcolorbox}
\usepackage{tikz}
\usetikzlibrary{arrows.meta,positioning,shapes.geometric,calc}
\usepackage[european]{circuitikz}
\usepackage{tikz-feynman}
\usepackage{pgfplots}
\pgfplotsset{compat=1.18}

% ---------------------------------------------------------------------
%  Farben
% ---------------------------------------------------------------------
\definecolor{dunkel}{HTML}{16213E}   % Hintergrund dunkler Folien, Titel
\definecolor{akzent}{HTML}{E9A23B}   % Akzent: Balken, Aufzählungszeichen
\definecolor{blau}{HTML}{3D7EA6}     % Zweitfarbe: Code, Hervorhebungen
\definecolor{papier}{HTML}{FBFBF8}   % Folienhintergrund
\definecolor{hell}{HTML}{EDF1F5}     % Kästen
\definecolor{grau}{HTML}{6B7280}     % Nebentext
\definecolor{codebg}{HTML}{F2F4F7}   % Code-Hintergrund

% ---------------------------------------------------------------------
%  Beamer-Layout
% ---------------------------------------------------------------------
\setbeamersize{text margin left=1.3cm, text margin right=1.3cm}
\setbeamercolor{background canvas}{bg=papier}
\setbeamercolor{normal text}{fg=dunkel}
\setbeamercolor{structure}{fg=dunkel}
\setbeamercolor{alerted text}{fg=akzent!80!black}
\setbeamerfont{frametitle}{size=\Large,series=\bfseries}
\setbeamertemplate{navigation symbols}{}

% Folientitel mit Akzentbalken
\setbeamertemplate{frametitle}{%
  \vspace*{0.7em}%
  {\usebeamerfont{frametitle}\color{dunkel}\insertframetitle\par}%
  \vspace{-0.1em}%
  {\color{akzent}\rule{2.4em}{2.5pt}}\par
  \vspace*{-0.6em}}

% Fußzeile mit Fortschrittsbalken
\setbeamertemplate{footline}{%
  \begin{tikzpicture}
    \useasboundingbox (0,0) rectangle (\paperwidth,0.8cm);
    \fill[hell] (0,0) rectangle (\paperwidth,2.5pt);
    \pgfmathsetmacro{\fortschritt}{min(1,\the\value{framenumber}/\inserttotalframenumber)}
    \fill[akzent] (0,0) rectangle (\fortschritt\paperwidth,2.5pt);
    \node[anchor=west, font=\tiny, text=grau] at (1.3cm,0.42cm) {Intro to \LaTeX};
    \node[anchor=east, font=\tiny, text=grau] at (\paperwidth-1.3cm,0.42cm)
      {\insertframenumber\,/\,\inserttotalframenumber};
  \end{tikzpicture}}

% Aufzählungszeichen
\setbeamertemplate{itemize item}{\color{akzent}\raisebox{0.15ex}{\rule{0.42em}{0.42em}}}
\setbeamertemplate{itemize subitem}{\color{grau}\raisebox{0.25ex}{\rule{0.3em}{0.3em}}}
\setbeamertemplate{enumerate item}{%
  \tikz[baseline=(n.base)]\node[fill=akzent, circle, inner sep=1.4pt, minimum size=1.25em,
    text=white, font=\scriptsize\bfseries] (n) {\insertenumlabel};}
\setlength{\leftmargini}{1.4em}

% Titelfolie
\setbeamertemplate{title page}{%
  \begin{tikzpicture}[remember picture, overlay]
    \fill[dunkel] (current page.south west) rectangle (current page.north east);
    % dezente Formeln im Hintergrund
    \begin{scope}[text=white, opacity=0.07]
      \node[font=\Huge, rotate=8] at ([xshift=-3.8cm,yshift=2.6cm]current page.east)
        {$i\hbar\,\partial_t \Psi = \hat{H}\Psi$};
      \node[font=\LARGE, rotate=-6] at ([xshift=-2.6cm,yshift=0.4cm]current page.east)
        {$\nabla \cdot \vec{E} = \rho / \varepsilon_0$};
      \node[font=\Huge, rotate=4] at ([xshift=-4.2cm,yshift=-1.7cm]current page.east)
        {$E = mc^2$};
      \node[font=\LARGE, rotate=-10] at ([xshift=-1.8cm,yshift=-3.0cm]current page.east)
        {$\oint \vec{B} \cdot \mathrm{d}\vec{s} = \mu_0 I$};
    \end{scope}
    \node[anchor=south west, inner sep=0, text=white,
          font=\fontsize{42}{46}\selectfont\bfseries] (t)
      at ([xshift=1.5cm,yshift=0.5cm]current page.west) {\inserttitle};
    \fill[akzent] ([yshift=-0.45cm]t.south west) rectangle ++(3cm,0.13cm);
    \node[anchor=north west, inner sep=0, text=white!85!dunkel, font=\Large]
      at ([yshift=-0.9cm]t.south west) {\insertsubtitle};
    \node[anchor=south west, inner sep=0, text=white!65!dunkel, font=\small, align=left]
      at ([xshift=1.5cm,yshift=1.2cm]current page.south west)
      {\textbf{\insertauthor}\\[2pt]\insertinstitute\\[2pt]\insertdate};
  \end{tikzpicture}}

% Abschnittsfolien
\newcommand{\zweistellig}[1]{\ifnum#1<10 0\fi\number#1}
\AtBeginSection[]{%
  \begin{frame}[plain,noframenumbering]
    \begin{tikzpicture}[remember picture, overlay]
      \fill[dunkel] (current page.south west) rectangle (current page.north east);
      \node[anchor=south west, inner sep=0, text=akzent,
            font=\fontsize{64}{64}\selectfont\bfseries] (n)
        at ([xshift=1.5cm,yshift=0.1cm]current page.west) {\zweistellig{\value{section}}};
      \node[anchor=north west, inner sep=0, text=white,
            font=\fontsize{30}{34}\selectfont\bfseries]
        at ([yshift=-0.55cm]n.south west) {\insertsectionhead};
    \end{tikzpicture}
  \end{frame}}

% ---------------------------------------------------------------------
%  Kästen
% ---------------------------------------------------------------------
\newtcolorbox{kasten}[1]{enhanced, halign=flush left, colback=hell, colframe=akzent, boxrule=0pt,
  leftrule=3pt, arc=3pt, left=8pt, right=8pt, top=4pt, bottom=6pt,
  title={#1}, colbacktitle=hell, coltitle=dunkel, fonttitle=\bfseries,
  titlerule=0pt, toptitle=5pt, bottomtitle=0pt}

\newtcolorbox{merke}{enhanced, halign=flush left, colback=akzent!14!papier, colframe=akzent,
  boxrule=0pt, leftrule=3pt, arc=3pt, left=10pt, top=5pt, bottom=5pt}

\newtcolorbox{karte}[3][akzent]{enhanced, halign=flush left, colback=white, colframe=hell,
  boxrule=0.8pt, arc=4pt, left=8pt, right=8pt, top=6pt, bottom=2pt,
  borderline north={3pt}{0pt}{#1}, fontupper=\small,
  before upper app={\setlength{\itemsep}{0pt}},
  before upper={{\large\bfseries #2}\par{\scriptsize\color{grau}#3}\par\vspace{0.2em}}}

\newtcolorbox{bildkarte}{enhanced, colback=white, colframe=hell, boxrule=0.8pt,
  arc=4pt, halign=center, valign=center}

% kleine Markierungen
\newcommand{\pill}[2][akzent]{%
  \tikz[baseline=(p.base)]\node[fill=#1, text=white, rounded corners=6pt,
    inner xsep=6pt, inner ysep=2.5pt, font=\scriptsize\bfseries] (p) {#2};}
\newcommand{\zeittag}[1]{%
  \begin{tikzpicture}[remember picture, overlay]
    \node[anchor=north east, fill=dunkel, text=white, rounded corners=6pt,
      inner xsep=7pt, inner ysep=3pt, font=\scriptsize\bfseries]
      at ([xshift=-1.3cm,yshift=-1.05cm]current page.north east) {#1};
  \end{tikzpicture}}

% Zeile im Ablauf-Zeitstrahl
\newcommand{\zeile}[5]{%
  \node[anchor=east, font=\small\bfseries, text=grau] at (1.9,#1) {#2};
  \fill[#5] (2.3,#1) circle (4.5pt);
  \node[anchor=west, font=\bfseries, text=#5!70!dunkel] at (2.8,#1) {#3};
  \node[anchor=west, font=\small, text=grau] at (7.6,#1) {#4};}

% ---------------------------------------------------------------------
%  Code-Darstellung
% ---------------------------------------------------------------------
\lstdefinestyle{latex}{
  language=[LaTeX]TeX,
  basicstyle=\ttfamily\footnotesize\color{dunkel},
  commentstyle=\color{grau}\itshape,
  texcsstyle=*\color{blau}\bfseries,
  keywordstyle=\color{blau}\bfseries,
  columns=fullflexible, keepspaces=true, breaklines=true,
  showstringspaces=false,
  moretexcs={documentclass,usepackage,section,subsection,label,ref,eqref,
             begin,end,item,frac,sqrt,int,sum,vec,partial,qty,unit,braket,
             ket,bra,newcommand,includegraphics,textbf,emph,maketitle,title,
             author,text,nabla,varepsilon,mathrm,hat},
  literate={ä}{{\"a}}1 {ö}{{\"o}}1 {ü}{{\"u}}1
           {Ä}{{\"A}}1 {Ö}{{\"O}}1 {Ü}{{\"U}}1 {ß}{{\ss}}1
}
\lstdefinestyle{log}{
  basicstyle=\ttfamily\footnotesize\color{dunkel},
  columns=fullflexible, keepspaces=true}

% Nur Code
\newtcblisting{code}[1][]{enhanced, listing only, colback=codebg,
  colframe=codebg, boxrule=0pt, arc=4pt, left=8pt, right=6pt, top=3pt,
  bottom=3pt, listing options={style=latex}, #1}

% Code links, gesetztes Ergebnis rechts
\newtcblisting{codeerg}[1][]{bicolor, listing side text, colback=codebg,
  colbacklower=white, colframe=codebg, boxrule=1pt, arc=4pt,
  lefthand ratio=0.57, sidebyside gap=14pt, sidebyside align=center seam,
  segmentation hidden, left=8pt, right=8pt, top=14pt, bottom=6pt,
  listing options={style=latex},
  overlay={%
    \node[anchor=north west, inner sep=5pt, font=\tiny\bfseries, text=grau]
      at (frame.north west) {CODE};
    \node[anchor=north west, inner sep=5pt, font=\tiny\bfseries, text=grau]
      at (frame.north west-|segmentation.east) {ERGEBNIS};},
  #1}

% ---------------------------------------------------------------------
%  Titelangaben
% ---------------------------------------------------------------------
\title{Intro to \LaTeX}
\subtitle{Workshop für Physik-Erstis}
\author{Gigi}
\institute{jDPG Regionalgruppe Tübingen}
\date{\today}
